\subsection{Uniform S-unit field norms and the exact capitulation obstruction}

\begin{statement}{E.79}{Proved}{All cubic and ninth-degree S-class norm quotients in the base-changed tower}
Use the actual fields $B_d=L_d(\xi)$ of E.76, with
$\xi^3=\zeta$, $\beta_{d+1}^3=\beta_d$, and $\beta_1=a\xi$,
$a^3=17$. Let $S_d$ be the six primes above $3$ and $17$, and put
\[
 \mathscr U_d=\mathcal O_{B_d}[1/3,1/17]^\times,
 \qquad A_d=\operatorname{Cl}_{S_d}(B_d)\{3\}.
 \tag{E79.1}
\]
For $r=1,2$, write $n=3^r$, $E=B_{d+r}$, $F=B_d$,
$G=\langle\sigma\rangle=\operatorname{Gal}(E/F)$, and
\[
 K_{d,r}=\mathscr U_d\cap N_{E/F}E^\times.
 \tag{E79.2}
\]
In every one of these actual cyclic extensions, the six Hilbert
characters, with the first argument in $\mathscr U_d$ and the
second argument $\beta_d$, map onto the entire product-one subgroup:
\[
 \rho_{d,r}(u)=\bigl((u,\beta_d)_{F_v,n}\bigr)_{v\in S_d},
 \qquad
 \rho_{d,r}(\mathscr U_d)
   =\ker\!\left(\mu_n^6\xrightarrow{\prod}\mu_n\right),
 \quad \ker\rho_{d,r}=K_{d,r}.
 \tag{E79.3}
\]
Consequently, $[\mathscr U_d:K_{d,r}]=n^5$ at every level.
The actual ideal norm induces the isomorphism
\[
 N_{E/F}: (A_{d+r})_G\xrightarrow{\ \sim\ } A_d,
 \qquad
 \ker(N:A_{d+r}\to A_d)=(\sigma-1)A_{d+r}.
 \tag{E79.4}
\]
In particular,
$|(A_{d+r})^G|=|A_d|$. At $d=1$, the two displayed
coinvariant groups are both $\mathbb Z/3$.

Here are complete finite coordinates for the norm calculation at
$B=B_1$. Use the certified ordinary-unit generators
$(\xi,u_1,\ldots,u_8)$ of E.76 and the six literal S-unit
generators $(s_1,\ldots,s_6)$ in the accompanying certificate.
Together with $-1$, they generate $\mathscr U_1$; their valuation
matrix, in the order of E.76, is
\[
 V=\begin{pmatrix}
 3&2&2&1&1&1\\0&1&0&0&0&0\\0&0&1&0&0&0\\
 0&0&0&1&0&0\\0&0&0&0&1&0\\0&0&0&0&0&1
 \end{pmatrix}.
 \tag{E79.5}
\]
The six full field-norm characters on these fifteen generators
are $[U_n\mid C_n]$ modulo $n$. The ordinary-unit blocks satisfy
$U_3=M_3$ and
$U_9=\operatorname{diag}(1,1,1,8,8,7)M_9$, with the $M_n$
of E.76. Thus the change of local quotient generators is explicit.
The new blocks are
\[
 C_3=\begin{pmatrix}
 1&1&0&0&0&2\\1&0&1&1&1&1\\1&2&2&0&0&1\\
 2&2&0&1&0&1\\2&2&0&0&1&1\\1&1&0&0&0&1
 \end{pmatrix},\qquad
 C_9=\begin{pmatrix}
 2&2&6&6&6&1\\8&3&5&8&2&8\\2&4&1&0&6&2\\
 1&1&6&2&0&5\\7&1&3&0&5&8\\4&7&3&0&0&1
 \end{pmatrix}.
 \tag{E79.6}
\]
The two image groups are $(\mathbb Z/n)^5$. Their full
fifteen-coordinate kernel lattices and all integral Smith maps
and inverses are retained. The allowed exponent vectors of the
$s_i$, after multiplication by an ordinary unit, are respectively
$K_n\mathbb Z^6$, where
\[
 K_3=\begin{pmatrix}
 1&0&0&0&0&0\\0&1&0&0&0&0\\0&0&1&0&0&0\\
 0&0&0&3&0&1\\0&0&0&0&3&1\\0&0&0&0&0&1
 \end{pmatrix},\quad
 K_9=\begin{pmatrix}
 1&0&0&0&0&0\\0&3&1&2&2&0\\0&0&1&0&0&0\\
 0&0&0&3&0&1\\0&0&0&0&3&1\\0&0&0&0&0&1
 \end{pmatrix}.
 \tag{E79.7}
\]
Hence $\mathbb Z^6/K_3\mathbb Z^6=(\mathbb Z/3)^2$
and $\mathbb Z^6/K_9\mathbb Z^6=\mathbb Z/9\oplus\mathbb Z/3$.
In the actual six valuation coordinates $w$, the full criterion is
\[
 \begin{aligned}
 w_1+w_2+w_3+2(w_4+w_5+w_6)&\equiv0\pmod3,\\
 w_4&\equiv w_5\equiv w_6\pmod3,\\
 3w_2-3w_3-2w_4-2w_5+4w_6&\equiv0\pmod9\quad(n=9).
 \end{aligned}
 \tag{E79.8}
\]
The first line is exactly the principal S-divisor condition.
For each $n$, the retained $9$ by $6$ matrix $R_n$ satisfies
$U_nR_n+C_nK_n=0$ modulo $n$. It gives the literal element
\[
 z\longmapsto
 \prod_{i=1}^{9}(\xi,u_1,\ldots,u_8)_i^{(R_nz)_i}
 \prod_{i=1}^{6}s_i^{(K_nz)_i}\in K_{1,r},
 \qquad \operatorname{div}_{S_1}=VK_nz.
 \tag{E79.9}
\]
The same included elements have the same field-norm characters
in every $B_{d+r}/B_d$.

An actual coherent prime chain makes (E79.4) concrete. In the
polynomial model $t=a-\xi$ of E.76, choose the degree-one prime
$\mathfrak p_1$ above $73$ with residues
\[
 t\mapsto33,\quad \xi\mapsto55,\quad a\mapsto15,
 \quad\beta_1\mapsto22,\qquad [\mathfrak p_1]=(2,1)
       \text{ in }\operatorname{Cl}(B).
 \tag{E79.10}
\]
For every $d$, extend this residue map by
$\beta_d\mapsto22$ for odd $d$ and $\beta_d\mapsto63$ for even
$d$. This defines a degree-one prime $\mathfrak p_d$ of $B_d$;
$N_{B_{d+1}/B_d}\mathfrak p_{d+1}=\mathfrak p_d$.
The three-primary components $[\mathfrak p_2]_3$ and
$[\mathfrak p_3]_3$ of their ideal classes generate
$A_2$ and $A_3$ as modules over
$\mathbb Z_3[\operatorname{Gal}(B_2/B)]$ and
$\mathbb Z_3[\operatorname{Gal}(B_3/B)]$, respectively. Their
annihilator ideals are not evaluated here. If $I_r$ is the
kernel of the actual surjection
$\mathbb Z_3[G]\to A_{1+r}$, $1\mapsto[\mathfrak p_{1+r}]_3$,
then the extension of the base generator is the explicit element
$j[\mathfrak p_1]_3=(1+\sigma+\cdots+\sigma^{n-1})
[\mathfrak p_{1+r}]_3$. Thus its remaining capitulation test is
membership of this norm polynomial in the actual ideal $I_r$.

The exact remaining obstruction to a norm preimage being an
S-unit is as follows. Let $j:A_d\to A_{d+r}$ extend ideals.
There is a surjective homomorphism, with its construction in the
proof,
\[
 \begin{split}
 \omega_{d,r}:K_{d,r}&\longrightarrow (A_{d+r})^G/jA_d,\\
 \ker\omega_{d,r}&=N_{E/F}\mathscr U_{d+r},\\
 \bigl|K_{d,r}/N\mathscr U_{d+r}\bigr|
   &=|\ker j|,\qquad
 [\mathscr U_d:N\mathscr U_{d+r}]=n^5|\ker j|.
 \end{split}
 \tag{E79.11}
\]
At $d=1$ the obstruction group has order dividing three.
In particular every element of $K_{1,r}$ has its cube in
$N\mathscr U_{1+r}$. One explicit family of field norms uses
$c_1=\beta_1+3\zeta+1$:
\[
 \xi c_1\in K_{d,1},\qquad
 \xi u_4u_5u_7^6u_8^6c_1\in K_{d,2}\qquad(d\ge1),
 \tag{E79.12}
\]
where these base elements are included in $B_d$.
The uncorrected $c_1$ has residue
$(1+3\zeta)^{96}=\zeta\ne1$ at every prime above $17$,
and therefore is not a norm in the cubic step. This has been
corrected in (E79.12) by an actual unit, before applying
(E79.11).

All finite-index lattice and obstruction sequences above remain
exact after completion for multiplication by $3^m$; the completed
S-unit groups are compact. The norm-coherent subgroup of E.69
still requires compatible actual S-unit norms. Its original
eventual exponent $h$, the higher full S-class groups, and the
capitulation maps $j$ remain unevaluated. The literal norm and
extension maps in E76.11 and the positional maps of E.46 and
E.72 are the comparisons to those original systems. The
manuscript anchor is LM~\S51.16, fols.~733v--736v,
scans~172--178.
\end{statement}

\begin{proof}
We first prove the uniform assertion with its exact local
normalization. The field $B_d$ contains $\mu_9$, and
$B_{d+r}=B_d(\beta_{d+r})$ with
$\beta_{d+r}^{\,n}=\beta_d$. At each of the three primes above
$17$, $\beta_d$ is a uniformizer, so this polynomial is
Eisenstein. At a prime above $3$, use
$(Y-1)^n-\beta_d$: the constant term $-1-\beta_d$ has valuation
one and each intermediate coefficient is divisible by $3$.
The uniformizer assertion for $\beta_d+1$ follows from the
Eisenstein polynomial over $\mathbb Q_3(\zeta)$ in E.76.
Thus every one of the six local extensions has degree $n$ and
is totally ramified. There are no other ramified primes,
since the discriminant of $Y^n-\beta_d$ is supported above
$3$ and $17$. All infinite places are complex. The inclusions
$B_1\subset B_d$ have, at these six places, the full local
degree $3^{d-1}$ and a single prime over each base prime.

Fix the arithmetic local Artin maps. Our Hilbert symbol is
\[
 (x,y)_{F_v,n}
  =\frac{\operatorname{Art}_{F_v}(x)(y^{1/n})}{y^{1/n}},
 \qquad (x,y)_{F_v,n}=(y,x)_{F_v,n}^{-1}.
 \tag{E79.13}
\]
For a finite extension of local fields $M/F_v$, reciprocity
gives the square
\[
 \begin{array}{ccc}
 M^\times&\xrightarrow{\operatorname{Art}_M}&G_M^{\mathrm{ab}}\\
 N_{M/F_v}\downarrow&&\downarrow\iota^{\mathrm{ab}}\\
 F_v^\times&\xrightarrow{\operatorname{Art}_{F_v}}&G_{F_v}^{\mathrm{ab}}.
 \end{array}
 \tag{E79.14}
\]
Here the right arrow is induced by inclusion of absolute
Galois groups. On a finite Galois overfield, this square
is precisely corestriction compatibility for the Tate
cup-product isomorphism
$\widehat H^{-2}(G,\mathbb Z)\to\widehat H^0(G,L^\times)$:
corestriction on $\widehat H^{-2}$ is the inclusion-induced
map on abelianizations, and on $\widehat H^0$ it is norm.
The fundamental class restricts to the fundamental class,
and the projection formula
$\operatorname{cor}(z\cup\operatorname{res}u)
 =\operatorname{cor}(z)\cup u$ proves the square.
These are the proved compatibility formulas immediately
before and after Milne's III Theorem~3.1; III Lemma~3.2
also states the square but does not print a separate proof
\cite{E_MilneCFT}. Passing to finite Galois overfields and
then their inverse limit gives (E79.14) without a Galois
assumption on $M/F_v$.

For $x\in F_v^\times$ and $y\in M^\times$, applying that
square to the root $x^{1/n}$ first proves
$(y,x)_{M,n}=(N_{M/F_v}y,x)_{F_v,n}$. Taking inverses in
(E79.13) therefore gives the second-argument formula
\[
 (x,y)_{M,n}=(x,N_{M/F_v}y)_{F_v,n}.
 \tag{E79.15}
\]
This fixes both arguments and the orientation. In the present
local tower the minimal polynomial of $\beta_d$ over $B_1$
is $Y^{3^{d-1}}-\beta_1$, and its odd degree gives the exact
norm $N\beta_d=\beta_1$. Hence, for every included
$u\in\mathscr U_1$ and each of the six compatible places,
\[
 (u,\beta_d)_{(B_d)_v,n}
      =(u,\beta_1)_{(B_1)_{v_1},n}.
 \tag{E79.16}
\]
The same chosen $\mu_n\subset B_1$ occurs on both sides.

We next verify the finite input to this identity. The full
unconditional class and unit certificate in E.76 gives
$\operatorname{Cl}(B)=(\mathbb Z/3)^2$ and the six
ramified-prime coordinates $(1,0)$ and $(2,0)$. Thus
$A_1=\mathbb Z/3$. The certified S-unit algorithm gives six
additional generators, whose actual ideal identities are
the columns of (E79.5). Their determinant is three, and their
image is exactly the kernel of
\[
 \mathbb Z^6\longrightarrow\operatorname{Cl}(B),\qquad
 w\longmapsto(w_1+w_2+w_3+2w_4+2w_5+2w_6,0)\pmod3.
 \tag{E79.17}
\]
Consequently these six generators, the complete ordinary
unit basis, and $-1$ generate all S-units. The source and
output include their literal polynomials and verify every
principal ideal identity \cite{E_PARICertifiedUnits}.

To extend the local unit characters of E.76 to the new
S-units, divide by an actual norm uniformizer before using
the recorded unit logarithm. At the wild places use
$\pi=\beta_1+1=N(Y+1)$; at the tame places use
$\beta_1=N(Y)$, where $Y^n=\beta_1$. Both identities have
positive sign because $n$ is odd. For a local input $z$,
the unit to be evaluated is $z/\pi^{v(z)}$, respectively
$z/\beta_1^{v(z)}$. These are exactly the evaluations in
(E79.6). The tame quotient generators in the ninth-degree
run differ from those of E.76 by the displayed diagonal
unit matrix, which the program checks on all nine ordinary
unit generators. No common-root identification is imposed
on the six computational rows.

The HNF lattices $n\mathbb Z^6+[U_n\mid C_n]\mathbb Z^{15}$
have index $n$ in $\mathbb Z^6$. Their image modulo $n$
therefore has order $n^5$; the complete Smith maps identify
it with $(\mathbb Z/n)^5$. Multiplying each row by the
unit that changes its local quotient generator to
(E79.13) preserves this order. Global reciprocity puts
the common-root image in the product-one subgroup, whose
order is $n^5$. Thus equality holds. Formula (E79.16)
includes this entire group in $\rho_{d,r}(\mathscr U_d)$
at every higher level, and the same reciprocity bound
proves (E79.3). At places outside $S_d$, an S-unit is a
local norm because the extension is unramified and norm
is onto the unit group. The cyclic Hasse norm theorem
then identifies the simultaneous kernel with (E79.2).
We use the exact statement and full proof in Milne VIII
Theorem~3.1; the same character-kernel argument is written
in Li--Ouyang--Xu--Zhang, Lemma~2.8
\cite{E_MilneCFT,E_LOXZKummer}.

For completeness, here is a direct class-field proof of
the actual norm isomorphism. Let $\mathcal H_E^S/E$ be
the maximal abelian extension unramified everywhere and
split at all primes of $S_{d+r}$. Let $\mathcal G/F$ be
its maximal subfield abelian over $F$, and set
$\Gamma=\operatorname{Gal}(\mathcal G/F)$. The extension
$\mathcal H_E^S/F$ is Galois. Its decomposition group at
any prime above $S_d$ maps isomorphically onto $G$,
because that prime is totally ramified in $E/F$ and
split in $\mathcal H_E^S/E$. This gives a splitting of
the finite group extension by $\operatorname{Cl}_S(E)$.
Its commutator subgroup is consequently
$(\sigma-1)\operatorname{Cl}_S(E)$, and
\[
 \ker(\Gamma\to G)=\bigl(\operatorname{Cl}_S(E)\bigr)_G.
 \tag{E79.18}
\]
Reciprocity describes $\Gamma$ as the quotient of finite
ideles of $F$ by
\[
 F^\times\!\left(\prod_{v\in S_d}N_{E_w/F_v}E_w^\times
                  \prod_{v\notin S_d}\mathcal O_{F_v}^\times\right).
 \tag{E79.19}
\]
This description can be checked in both directions. Its
class field contains $E$. At $v\in S_d$, its local
decomposition group has order at most $n$ and maps onto
the degree-$n$ local extension $E_w/F_v$, so its local
extension over $E$ is trivial. Outside $S_d$ it is
unramified. Conversely every abelian subextension of
$\mathcal H_E^S/F$ kills the local norms in (E79.19),
by (E79.14), and kills the unit groups outside $S_d$.
This proves the asserted maximality and the actual
reciprocity quotient.

Mapping a finite idele to its divisor outside $S_d$
gives an exact sequence
\[
 \mathscr U_d\longrightarrow
 \prod_{v\in S_d}F_v^\times/N E_w^\times
 \longrightarrow\Gamma\longrightarrow
 \operatorname{Cl}_{S_d}(F)\longrightarrow0.
 \tag{E79.20}
\]
The product has order $n^6$, and (E79.3) says that the
first arrow has image order $n^5$. Thus the kernel of
the last arrow has order $n$. The local Artin image
at any one prime in $S_d$ maps onto $G$; it is contained
in this kernel, so that the kernel maps isomorphically
to $G$. It follows that the restriction
$\ker(\Gamma\to G)\to\operatorname{Cl}_{S_d}(F)$ is
an isomorphism. For a prime $w$ of $E$ outside $S$,
its Frobenius restricts to the $f(w/v)$th power of the
Frobenius at $v$; its image in the last class group is
$f(w/v)[v]=N[w]$. Thus this is the actual ideal norm
on (E79.18). Taking three-primary parts proves
(E79.4). Equivalently, norm surjectivity follows by
base changing the S-split Hilbert class field of $F$:
its intersection with $E$ is $F$ because every nontrivial
subextension of $E/F$ ramifies totally at $S_d$. The
resulting S-split unramified extension of $E$ has
restriction map on Frobenii equal to ideal norm.
For a finite group and an endomorphism $\sigma-1$,
kernel and cokernel have equal order. This proves the
invariant-order assertion as well. It agrees with
Gras' formula in \cite[Theorem~2.6]{E_LOXZKummer},
with $D$ generated by the six ramified primes; here
$\Lambda_D=\mathscr U_d$ and the actual norm of $D$
is the full base S-divisor group.

We give the full lattice argument behind (E79.7)--(E79.9).
Put $H_n=n\mathbb Z^6+U_n\mathbb Z^9$. A vector
$t\in\mathbb Z^6$ allows an ordinary-unit correction
precisely when $C_nt\in H_n$. The retained integer
kernel computation gives
$\{t:C_nt\in H_n\}=K_n\mathbb Z^6$.
Both inclusions are certified: $H_n^{-1}C_nK_n$ is
integral, and
\[
 |\det K_n|=
 \frac{[\mathbb Z^6:H_n]}
 {[\mathbb Z^6:H_n+C_n\mathbb Z^6]}.
 \tag{E79.21}
\]
The indices are $9$ and $27$. Integral Smith
transformations and their inverses give exactly the
quotients stated after (E79.7); they also give
$\mathbb Z^6/VK_n\mathbb Z^6=(\mathbb Z/3)^3$
for $n=3$, and $\mathbb Z/9\oplus(\mathbb Z/3)^2$
for $n=9$. Applying $V^{-1}$ and then $K_n^{-1}$ gives
(E79.8), including its integer divisibility conditions.
The computed $R_n$ gives (E79.9) directly. The complete
fifteen-coordinate kernel has determinant $n^5$, so it
contains every S-unit field norm, with the relation
$\xi^9=1$ already lying in that kernel. The matrix
identities and all their inverse coordinate maps are
independently checked by \texttt{e\_sclass\_norms.py}.

The prime certificate in (E79.10) is an actual ideal
of norm $73$ in the retained integral basis. Reducing
that basis gives its quotient map to $\mathbb F_{73}$;
the ideal is exactly the kernel, as follows from its
determinant $73$. The certified class coordinates are
$(2,1)$, which maps to a generator of $A_1$.
The two identities $63^3=22$ and $22^3=63$ in
$\mathbb F_{73}$ prove compatibility of all the
displayed roots. Their order is eight. At $73$ the
polynomials $Y^{3^{d-1}}-\beta_1$ are etale, so these
residue maps define actual degree-one primes of the
maximal orders. Ideal norms send each to its predecessor.
At the first two levels, (E79.4) shows that the three-primary
component of the prime class generates the coinvariants.
Projection to that component is the canonical projection
in the primary decomposition of the finite ideal class
group, and commutes with ideal norm and extension.
To see the asserted
cyclic-module generation without an extra assumption,
let $Q$ be its quotient in the finite three-group
$A_{1+r}$. It satisfies
$Q=(3,\sigma-1)Q$. Modulo $3$ the augmentation ideal
is nilpotent since $(\sigma-1)^n=0$ in
$\mathbb F_3[G]$. Iterating therefore gives $Q=3Q$,
and finiteness gives $Q=0$. This proves the stated
surjection from $\mathbb Z_3[G]$ with the literal
prime component as image of one. The base prime splits
completely in both extensions, since its displayed root
and all its $n$th-root-of-unity multiples lie in
$\mathbb F_{73}$. Extension of its ideal is therefore
the product of these $n$ conjugate upper primes. This
proves the stated norm-polynomial formula for $j$.

We finally prove the norm-defect and capitulation maps.
Let $\operatorname{Div}_{E,S}$ denote divisors outside
$S_{d+r}$. On each orbit of unramified primes, norm
is residue degree times the sum of coefficients.
Therefore
\[
 \ker\bigl(N:\operatorname{Div}_{E,S}
                 \to\operatorname{Div}_{F,S}\bigr)
       =(\sigma-1)\operatorname{Div}_{E,S},\qquad
 (\operatorname{Div}_{E,S})^G=j\operatorname{Div}_{F,S}.
 \tag{E79.22}
\]
The second identity uses exactly that all ramification
lies in $S$: extension of a base prime is the sum of
its conjugate upper primes, each with coefficient one.
For $u\in K_{d,r}$ choose $x\in E^\times$ with $Nx=u$.
Then $\operatorname{div}_{E,S}(x)$ has norm zero, so
choose $D$ with
$\operatorname{div}_{E,S}(x)=(\sigma-1)D$.
The class $[D]$ is $G$-invariant. Its image modulo
$j\operatorname{Cl}_{S_d}(F)$ is independent of $D$
by the second identity in (E79.22). If $x$ changes
to another norm preimage, Hilbert's theorem~90 writes
the quotient as $\sigma(z)/z$. Replacing $D$ by
$D+\operatorname{div}_{E,S}(z)$ shows independence
of this choice too. This defines $\omega$.

Every invariant ideal class has a representative $D$
with $(\sigma-1)D=\operatorname{div}_{E,S}(x)$ for
some $x$, and then $Nx$ is an S-unit; thus $\omega$
is onto. Its vanishing says that
$D=\operatorname{div}_{E,S}(z)+jD_0$. The element
$x/(\sigma z/z)$ then has divisor zero outside $S$
and the same norm $u$, proving the kernel formula.
Conversely an S-unit norm has the zero representative
$D=0$. The quotient is killed by $n$ because every
$u\in\mathscr U_d$ has $u^n=N(u)$ after inclusion
in $E$. Thus the quotient of full class groups just
constructed is its three-primary part, giving
exactly the groups in (E79.11).

There is a second exact map which records capitulation.
For a class $[D_0]\in\ker j$, choose $x$ with
$\operatorname{div}_{E,S}(x)=jD_0$ and send it to
the cocycle $g\mapsto g(x)/x$ in $\mathscr U_{d+r}$.
Principal changes and changing $x$ by an S-unit change
this by a coboundary. Hilbert~90 and the invariant
divisor identity in (E79.22) give the inverse. Hence
\[
 \ker(j:A_d\to A_{d+r})
       \simeq H^1(G,\mathscr U_{d+r}).
 \tag{E79.23}
\]
All these groups are three-primary since $Nj=n$.
The S-unit Herbrand quotient is $n^5$: in the
logarithmic permutation representation, the six
finite S-places contribute six trivial lattices,
the complex places contribute regular lattices,
and the product formula removes one trivial lattice.
The resulting quotient is $n^6/n$; finite-index
lattice changes and torsion have Herbrand quotient
one. This is also Milne VII Proposition~3.1, whose
full lattice proof applies with these exact local
degrees \cite{E_MilneCFT}. Therefore
\[
 \frac{[\mathscr U_d:N\mathscr U_{d+r}]}
 {|\ker j|}=n^5.
 \tag{E79.24}
\]
Together with (E79.3) this proves the norm-defect
size in (E79.11). It also follows directly from
$|(A_{d+r})^G|=|A_d|$ and the actual image of $j$.
At $d=1$, $A_1=\mathbb Z/3$, so the defect has
order dividing three and every cube assertion follows.

The literal corrections in (E79.12) are checked on
all six local characters; (E79.16) transports them
to every higher step. At the tame residue field
$\mathbb F_{17}[\zeta]/(\zeta^2+\zeta+1)$, direct
exponentiation gives $(1+3\zeta)^{96}=\zeta$.
Norm on residues in the totally ramified cubic
extension is cubing. This proves the asserted
failure for the uncorrected $c_1$. Multiplication
by $\xi$ cancels that residue and the three wild
characters as well, as the complete matrix check
records. The ninth-degree correction uses the same
certified ordinary-unit basis.

For the completion assertion, all S-unit groups and
all their displayed subgroups are finitely generated
abelian groups. The kernels of the finite characters
have finite index. Also $N\mathscr U_{d+r}$ contains
$\mathscr U_d^{\,n}$, so it has finite index in
$\mathscr U_d$. Completing these exact sequences
therefore retains their kernels, images and finite
three-primary quotients; explicitly the groups are
their finite three-torsion summands plus finite
products of $\mathbb Z_3$. In particular the
completed version of (E79.11) gives the same
obstruction group and the same zero criterion.
This proves an exact membership theorem for each
step. The all-level compatible norm condition
remains the actual inverse system of these maps,
with the capitulation maps and upper annihilator
ideals retained as unevaluated arithmetic data.
\end{proof}

